\chapter{Spatial Updated-Lagrangian Constitutive Interface}
\label{ch:spatial-updated-lagrangian-interface}

\section{Motivation}

Chapter~\ref{ch:continuum-constitutive-interface} introduced a continuum-side
carrier that exposed $\mathbf{F}$, $\det\mathbf{F}$ and explicit strain/stress
measure tags at the erased \texttt{Material<>} boundary. The next missing step
was to let the \texttt{UpdatedLagrangian} formulation use a genuinely spatial
pair instead of silently reusing the reference pair
$(\mathbf{E}, \mathbf{S})$.

For many future models that is not optional. Current-configuration
constitutive laws are the natural setting for:
\begin{itemize}
  \item large-displacement plasticity written in terms of Cauchy or Kirchhoff
        stress;
  \item rate-based models and objective stress rates;
  \item contact, FSI and other couplings that consume spatial tractions and
        current normals directly.
\end{itemize}

\section{Reference and Spatial Conjugate Pairs}

For the total-Lagrangian formulation the standard conjugate pair is
\[
  (\mathbf{E}, \mathbf{S}),
\]
where
\begin{align}
  \mathbf{F} &= \mathbf{I} + \nabla_0 \mathbf{u}, \\
  \mathbf{E} &= \frac{1}{2}\left(\mathbf{F}^T \mathbf{F} - \mathbf{I}\right).
\end{align}

For the updated-Lagrangian formulation the natural spatial pair is
\[
  (\mathbf{e}, \boldsymbol{\sigma}),
\]
where $\mathbf{e}$ is the Almansi strain,
\begin{align}
  \mathbf{b} &= \mathbf{F}\mathbf{F}^T, \\
  \mathbf{e} &= \frac{1}{2}\left(\mathbf{I} - \mathbf{b}^{-1}\right),
\end{align}
and $\boldsymbol{\sigma}$ is the Cauchy stress.

The two pairs are linked through Piola transformations:
\begin{align}
  \boldsymbol{\sigma} &= \frac{1}{J}\mathbf{F}\mathbf{S}\mathbf{F}^T, \\
  \mathbf{S} &= J\mathbf{F}^{-1}\boldsymbol{\sigma}\mathbf{F}^{-T},
\end{align}
with $J = \det \mathbf{F}$.

\section{Why the Spatial Carrier Matters}

The previous implementation already contained the mathematics needed to build
the spatial path, but the erased constitutive interface still compressed the
interaction into ``some Voigt strain'' leading to two semantic problems:
\begin{enumerate}
  \item the constitutive law could not tell whether the incoming stress measure
        was meant to be $\mathbf{S}$ or $\boldsymbol{\sigma}$;
  \item the updated-Lagrangian element path still assembled with a
        reference-style carrier at the element/material interface.
\end{enumerate}

This refactor resolves that ambiguity by making the active pair explicit at
compile time in the kinematic policy and explicit at runtime in the continuum
carrier tags.

\section{Updated-Lagrangian Carrier}

The continuum carrier now stores all of the following simultaneously:
\begin{itemize}
  \item engineering Voigt form of the active strain measure;
  \item infinitesimal strain;
  \item Green--Lagrange strain;
  \item Almansi strain;
  \item deformation gradient $\mathbf{F}$ and $J=\det\mathbf{F}$;
  \item the active strain and stress measure tags.
\end{itemize}

For \texttt{UpdatedLagrangian}, the active tags are
\[
  \texttt{StrainMeasureKind::almansi}, \qquad
  \texttt{StressMeasureKind::cauchy}.
\]

This means that a constitutive relation receiving the carrier can decide at the
entry point whether it should produce:
\begin{itemize}
  \item the reference response $\mathbf{S}$ and $\mathbb{C}$, or
  \item the spatial response $\boldsymbol{\sigma}$ and $\mathcal{c}$.
\end{itemize}

\section{Push-Forward of the Constitutive Tangent}

If a hyperelastic model is naturally written in terms of
$\mathbf{S} = \partial W / \partial \mathbf{E}$ and
$\mathbb{C} = \partial^2 W / \partial \mathbf{E}^2$, then the spatial tangent
used by the updated-Lagrangian assembly is the push-forward of
$\mathbb{C}$:
\[
  \mathcal{c}_{ijkl}
  = \frac{1}{J}
    F_{iI}F_{jJ}\,\mathbb{C}_{IJKL}\,F_{kK}F_{lL}.
\]

In the present implementation that transformation is handled at the relation
level for direct continuum-aware laws such as \texttt{HyperelasticRelation}.
The same direct overload that previously returned $\mathbf{S}$ now returns
$\boldsymbol{\sigma}$ when the carrier announces a Cauchy-conjugate path.

\section{Element-Level Updated-Lagrangian Assembly}

The updated-Lagrangian element route now uses:
\begin{align}
  \nabla_x N_I &= \nabla_X N_I \mathbf{F}^{-1}, \\
  \mathbf{b}(\nabla_x N) &= \text{spatial small-strain } B\text{-operator}, \\
  \delta W_{\text{int}}
    &= \int_{\Omega_n} \boldsymbol{\sigma} : \delta\mathbf{d}\, \dd v.
\end{align}

Because the implementation still integrates over the reference quadrature
domain attached to the original geometry, the conversion
\[
  \dd v = J \, \dd V
\]
is applied explicitly inside \texttt{ContinuumElement}. This is the reason the
element path now multiplies the updated-Lagrangian contributions by
$J = \det\mathbf{F}$.

At the discrete level, the internal force and tangent take the form
\begin{align}
  \mathbf{f}_{\text{int}}
    &= \int_{\Omega_0} J\, \mathbf{b}^T \tilde{\boldsymbol{\sigma}} \, \dd V, \\
  \mathbf{K}_{\text{mat}}
    &= \int_{\Omega_0} J\, \mathbf{b}^T \tilde{\mathcal{c}} \, \mathbf{b} \, \dd V, \\
  \mathbf{K}_{\sigma}
    &= \int_{\Omega_0} J\, \mathbf{k}_{\sigma}(\nabla_x N, \boldsymbol{\sigma}) \, \dd V.
\end{align}

\section{Architectural Consequences}

This step is important because it removes the last semantic shortcut in the
continuum large-displacement path:
\begin{quote}
\texttt{UpdatedLagrangian} is no longer ``TL in disguise'' at the generic
element/material boundary for hyperelastic continuum laws.
\end{quote}

The resulting architecture has three useful properties:
\begin{enumerate}
  \item \textbf{Static specialization is preserved.}
        The kinematic policy still selects the path at compile time.
  \item \textbf{The erased material interface remains additive.}
        Old laws keep working through the legacy \texttt{KinematicT} fallback.
  \item \textbf{Spatial finite-strain laws now have a natural slot.}
        Future Cauchy-based plasticity or rate-based models can consume the
        same carrier without pretending to be reference-measure laws.
\end{enumerate}

\section{Tests Added}

The new regression coverage verifies:
\begin{itemize}
  \item the updated-Lagrangian carrier exposes Almansi strain and Cauchy stress
        tags correctly;
  \item \texttt{HyperelasticRelation} returns the push-forwarded Cauchy stress
        and spatial tangent when the carrier requests the spatial path;
  \item the erased \texttt{Material<ThreeDimensionalMaterial>} handle forwards
        that spatial path without losing the richer semantics;
  \item \texttt{UpdatedLagrangian::evaluate\_from\_gradients()} now returns the
        spatial $B$ operator and the Almansi engineering strain.
\end{itemize}

\section{Remaining Limits}

The spatial route is now available for continuum hyperelasticity, but it is
not yet the full finite-strain inelastic infrastructure. The next steps remain:
\begin{itemize}
  \item a finite-strain \texttt{ConstitutiveState} split for continuum plasticity
        and damage;
  \item local residual/Jacobian contracts for nonlinear constitutive updates in
        the spatial setting;
  \item objective-rate or multiplicative-decomposition models when the library
        moves beyond hyperelastic and small-strain-inelastic ingredients.
\end{itemize}

\section{References}

\begin{thebibliography}{9}

\bibitem{bonet2008}
J.~Bonet, A.~J. Gil, and R.~D. Wood,
\emph{Nonlinear Solid Mechanics for Finite Element Analysis: Statics},
2nd ed., Cambridge University Press, 2008.

\bibitem{belytschko2014}
T.~Belytschko, W.~K. Liu, B.~Moran, and K.~Elkhodary,
\emph{Nonlinear Finite Elements for Continua and Structures},
2nd ed., Wiley, 2014.

\bibitem{simo1998}
J.~C. Simo and T.~J.~R. Hughes,
\emph{Computational Inelasticity},
Springer, 1998.

\end{thebibliography}
